\section{Related Work}
\label{sec:related_work}

Our work sits at the intersection of three research threads: transfer learning with frozen features, supervised dimensionality reduction, and the recent literature on efficient inference. We review each in turn.

\subsection{Transfer Learning and Linear Probing}

Using pretrained networks as fixed feature extractors was among the earliest demonstrations of deep transfer learning. Donahue et al.~\cite{donahue2014decaf} showed that features from a network trained on ImageNet generalize across visual recognition tasks. Sharif Razavian et al.~\cite{sharif2014cnn} extended this observation systematically, establishing that ``CNN features off-the-shelf'' are strong baselines for a range of tasks. Kornblith et al.~\cite{kornblith2019better} later studied the correlation between ImageNet accuracy and transfer performance, using logistic regression on frozen features as their evaluation protocol. Ericsson et al.~\cite{ericsson2021well} benchmarked self-supervised representations using the same linear probing methodology.

A notable gap in this literature is the treatment of feature dimensionality. All of the studies above feed the full feature vector to the linear classifier. To the best of our knowledge, none systematically evaluates whether intermediate dimensionality reduction improves the downstream task. The few that mention dimensionality do so in passing---Kornblith et al. observe that higher-dimensional models tend to transfer better, but do not test whether \emph{reducing} dimensions of a fixed model's features can be beneficial.

\subsection{Linear Discriminant Analysis}

LDA, introduced by Fisher~\cite{fisher1936lda} and extended to multiple classes by Rao~\cite{rao1948utilization}, remains one of the most widely used supervised dimensionality reduction methods. It seeks a linear projection that maximizes the ratio of between-class scatter $\mathbf{S}_b$ to within-class scatter $\mathbf{S}_w$:
\begin{equation}
\mathbf{W}^* = \argmax_{\mathbf{W}} \frac{|\mathbf{W}^\top \mathbf{S}_b \mathbf{W}|}{|\mathbf{W}^\top \mathbf{S}_w \mathbf{W}|}.
\label{eq:lda_objective}
\end{equation}
The solution consists of the leading eigenvectors of $\mathbf{S}_w^{-1}\mathbf{S}_b$, yielding at most $C-1$ discriminant directions for $C$ classes.

LDA's well-known limitations include the Gaussian class-conditional assumption, sensitivity to singular or ill-conditioned scatter matrices, and the hard cap of $C-1$ components. Several extensions address these issues. Regularized LDA~\cite{friedman1989regularized} replaces $\mathbf{S}_w$ with a shrinkage estimate $(1-\lambda)\mathbf{S}_w + \lambda\mathbf{I}$ to stabilize inversion. Penalized Discriminant Analysis~\cite{hastie1995penalized} imposes smoothness constraints. Kernel LDA~\cite{mika1999fisher} handles nonlinear class boundaries. Heteroscedastic extensions~\cite{loog2004lda_review} relax the equal-covariance assumption. Recent work by Li et al.~\cite{li2022survey_dr} surveys these and other variants.

In the context of deep features, LDA has been used primarily during training---for example, as a regularizer in loss functions~\cite{dorfer2016deep_lda} or in metric learning objectives~\cite{wen2016discriminative}. The simple approach of applying LDA \emph{after} feature extraction, as a post-hoc dimensionality reduction step, has received remarkably little attention in the modern deep learning era.

\subsection{Local Fisher Discriminant Analysis}

LFDA, proposed by Sugiyama~\cite{sugiyama2007lfda}, addresses the multimodal limitation of LDA by incorporating local structure. It uses an affinity matrix derived from nearest-neighbor distances to weight the scatter matrices, preserving the local geometry of each class. While theoretically appealing, LFDA introduces a hyperparameter (the number of neighbors) and incurs higher computational cost due to the $k$-NN graph construction. Our experiments show that this additional complexity does not translate into accuracy gains on frozen CNN features (Section~\ref{sec:experiments}).

\subsection{Neighbourhood Components Analysis}

NCA~\cite{goldberger2004nca} takes a different approach: it learns a linear transformation that maximizes a differentiable approximation to leave-one-out $k$-NN accuracy. This makes NCA a metric learning method rather than a projection-based one. While NCA can discover task-specific feature transformations, it is notoriously slow to train---in our experiments, NCA is 10--25$\times$ slower than LDA for negligible or negative accuracy differences.

\subsection{PCA as a Baseline}

PCA~\cite{pearson1901pca} is the standard unsupervised baseline for dimensionality reduction. It finds orthogonal directions of maximum variance, discarding low-variance dimensions. In the transfer learning setting, PCA is sometimes used as a preprocessing step~\cite{sharif2014cnn} but is rarely compared systematically against supervised alternatives like LDA. An important question that our study addresses is whether the label information used by LDA provides a meaningful advantage over PCA's purely variance-based criterion.

\subsection{Efficient Inference and Model Compression}

Our work is tangentially related to the literature on efficient inference. Knowledge distillation~\cite{hinton2015distilling}, pruning~\cite{han2015learning}, and quantization~\cite{jacob2018quantization} reduce the cost of running the backbone itself. Our approach is complementary: we reduce the cost of the \emph{downstream} classifier by operating in a lower-dimensional feature space. This is particularly relevant in settings where the backbone is a fixed, deployed model and only the classification head can be modified---a common scenario in edge deployment and multi-task systems.
